% unidades/28-voz-pasiva.tex
\chapter{📢 受身形 — Voz Pasiva}
\unidadheader{28}{受身形}{Voz Pasiva}{Hablar de acciones recibidas o hechos históricos}

\begin{tcolorbox}[vocabbox, title={📌 Vocabulario Nuevo}]
\begin{tabularx}{\linewidth}{llX}
\toprule
\textbf{Japonés} & \textbf{Español} & \textbf{Tipo} \\
\midrule
\vocabitem{叱る}{しかる}{regañar}{v. gr.1}
\vocabitem{褒める}{ほめる}{elogiar / alabar}{v. gr.2}
\vocabitem{頼む}{たのむ}{pedir un favor / encargar}{v. gr.1}
\vocabitem{誘う}{さそう}{invitar}{v. gr.1}
\vocabitem{踏む}{ふむ}{pisar}{v. gr.1}
\vocabitem{壊す}{こわす}{romper / destruir}{v. gr.1}
\vocabitem{盗む}{ぬすむ}{robar}{v. gr.1}
\vocabitem{発明する}{はつめいする}{inventar}{v. irr.}
\vocabitem{発見する}{はっけんする}{descubrir}{v. irr.}
\vocabitem{行われる}{おこなわれる}{llevarse a cabo}{v. gr.2}
\bottomrule
\end{tabularx}
\end{tcolorbox}

\subsection*{🗣️ Frases Clave}

\frase{\ruby{先生}{せんせい}に\ruby{褒}{ほ}められました。}{Fui elogiado por el profesor.}
\frase{\ruby{母}{はは}に\ruby{叱}{しか}られました。}{Fui regañado por mi madre.}
\frase{電車で\ruby{足}{あし}を\ruby{踏}{ふ}まれました。}{Me pisaron el pie en el tren.}
\frase{この\ruby{本}{ほん}は\ruby{多}{おお}くの\ruby{人}{ひと}に\ruby{読}{よ}まれています。}
  {Este libro es leído por mucha gente.}

\subsection*{⚙️ Gramática}

\patron{Forma pasiva: ser ~ado}
  {V (pasiva) + ます}
  {Grupo 1: う→あ段 + れる (ej: 書く→書かれる) \\
   Grupo 2: る→られる (ej: 食べる→食べられる) \\
   Grupo 3: する→される / くる→こられる}

\patron{Estructura de la oración pasiva}
  {受動者(receptor) は 能動者(agente) に V-れる}
  {\ej{\ruby{弟}{おとうと}は\ruby{兄}{あに}に\ruby{殴}{なぐ}られました。}
    {El hermano menor fue golpeado por el mayor.}
  \ej{この\ruby{寺}{てら}は600\ruby{年}{ねん}\ruby{前}{まえ}に\ruby{建}{た}てられました。}
    {Este templo fue construido hace 600 años.}}

\begin{tcolorbox}[ejerciciobox, title={📝 Ejercicios Rápidos}]
\begin{enumerate}
  \item Convierte a pasiva: \ruby{誘}{さそう} / \ruby{盗}{ぬす}む / \ruby{褒}{ほ}める.
  \item Di "Mi perro se comió mi tarea" en voz pasiva (lit. "Mi tarea fue comida por el perro").
  \item Traduce: この\ruby{ビル}{びる}は1990\ruby{年}{ねん}に\ruby{建}{た}てられました。
\end{enumerate}
\vspace{0.4em}
\textbf{Respuestas:}
1) \ruby{誘}{さそ}われる / \ruby{盗}{ぬす}まれる / \ruby{褒}{ほ}められる\quad
2) \ruby{宿題}{しゅくだい}は\ruby{犬}{いぬ}に\ruby{食}{た}べられました。\quad
3) Este edificio fue construido en 1990.
\end{tcolorbox}

\begin{tcolorbox}[kanjibox, title={🀄 Kanjis de la Unidad}]
\begin{tabularx}{\linewidth}{cllXX}
\toprule
\textbf{Kanji} & \textbf{On} & \textbf{Kun} & \textbf{Significado} & \textbf{Vocab} \\
\midrule
\kanjirow{受}{じゅ}{う}{recibir}{\ruby{受}{う}ける / \ruby{受身}{うけみ}}
\kanjirow{身}{しん}{み}{cuerpo / uno mismo}{\ruby{受身}{うけみ} / \ruby{出身}{しゅっしん}}
\kanjirow{盗}{とう}{ぬす}{robar}{\ruby{盗}{ぬす}む / \ruby{強盗}{ごうとう}}
\kanjirow{壊}{かい}{こわ}{romper}{\ruby{壊}{こわ}す / \ruby{破壊}{はかい}}
\kanjirow{建}{けん}{た}{construir}{\ruby{建}{た}てる / \ruby{建物}{たてもの}}
\bottomrule
\end{tabularx}
\end{tcolorbox}
